\section{Transferência de órbita}
\label{sec:hohmann}

No contexto de \emph{rendezvous and docking/berthing} (RVD/B), a transferência de Hohmann é utilizada como manobra inicial para que o \emph{chaser} ajuste sua trajetória orbital. 
Considera-se que o \emph{chaser} parte de uma órbita circular mais baixa e realiza a aproximação ao \emph{target} pela direção radial (\emph{R-bar}), isto é, por baixo em relação à referência LVLH.

Seja uma órbita inicial circular de raio $r_1$ e uma órbita final circular de raio $r_2$. 
Após o primeiro impulso ($\Delta v_1$), a espaçonave entra em uma órbita elíptica de transferência, cujo perigeu está em $r_1$ e apogeu em $r_2$. 
No instante em que atinge o apogeu, aplica-se o segundo impulso ($\Delta v_2$), circularizando a órbita em $r_2$. A velocidade na órbita circular inicial é dada por:
\begin{equation}
    v_1 = \sqrt{\frac{\mu}{r_1}},
\end{equation}

onde $\mu$ é o parâmetro gravitacional da Terra. 
A velocidade no perigeu da órbita de transferência é:
\begin{equation}
    v_{T,1} = \sqrt{2\mu \left(\frac{1}{r_1} - \frac{1}{2a_T}\right)},
\end{equation}
em que $a_T = \tfrac{1}{2}(r_1+r_2)$ é o semi-eixo maior da órbita de transferência.

A variação de velocidade no primeiro impulso é:
\begin{equation}
    \Delta v_1 = v_{T,1} - v_1.
\end{equation}

De forma análoga, a velocidade na órbita circular final é:
\begin{equation}
    v_2 = \sqrt{\frac{\mu}{r_2}},
\end{equation}
e a velocidade no apogeu da órbita de transferência é:
\begin{equation}
    v_{T,2} = \sqrt{2\mu \left(\frac{1}{r_2} - \frac{1}{2a_T}\right)}.
\end{equation}

A variação de velocidade no segundo impulso é:
\begin{equation}
    \Delta v_2 = v_2 - v_{T,2}.
\end{equation}

O tempo de voo da manobra de Hohmann corresponde a metade do período orbital da órbita de transferência:
\begin{equation}
    t_T = \pi \sqrt{\frac{a_T^3}{\mu}}.
\end{equation}

A variação total de velocidade requerida é:
\begin{equation}
    \Delta v_{H} = \Delta v_1 + \Delta v_2.
\end{equation}

Concluída a manobra de transferência de Hohmann, o \emph{chaser} encontra-se em uma órbita coplanar e ligeiramente abaixo da órbita do \emph{target}. 
A partir desse ponto, ajustes de fase são empregados para sincronizar angularmente as duas espaçonaves. 
Uma vez nesta posição, aplicam-se as técnicas de aproximação entre as espaçonaves e são descritas em termos do movimento relativo pelas equações de Hill/Clohessy-Wiltshire.

